
\documentclass[12pt]{article}
\usepackage{amsmath, amssymb}
\usepackage{geometry}
\geometry{margin=1in}
\title{SAT.O5 – Emergent Gauge Couplings from Filament Topology}
\author{SATO.4D Coherence Project}
\date{June 2025}

\begin{document}
\maketitle

\section*{1. Foundational Assumptions}

\begin{itemize}
    \item Gauge interactions are not fundamental but emerge from topological statistics of 1D filaments embedded in 4D spacetime.
    \item No gauge fields \( A_\mu \) or couplings \( g \) are inserted by hand.
    \item Linking and winding densities of filament bundles determine observable coupling constants.
\end{itemize}

\section*{2. Filament Topological Structures}

We define the following local topological densities:
\begin{itemize}
    \item \( \rho_{\text{wind}}(x) \): Winding number density (loops around compact directions)
    \item \( \rho_{\text{link}}(x) \): Pairwise linking density
    \item \( \rho_{\text{triplet}}(x) \): Triple-linking (Borromean) density
\end{itemize}

Each density reflects the statistical frequency of specific topological configurations within a volume element at \( x \in M \).

\section*{3. Dimensional and Structural Scaling}

Let:
\begin{itemize}
    \item \( \ell_f = \left( \frac{2A}{T} \right)^{1/3} \): Filament transverse scale
    \item \( \epsilon^2_f = \ell_f^2 \): Effective interaction area
\end{itemize}

Then the dimensionless gauge couplings are given by:

\[
g_{U(1)} \sim \frac{1}{\sqrt{\rho_{\text{wind}} \, \epsilon_f^2}}, \quad
g_{SU(2)} \sim \frac{1}{\sqrt{\rho_{\text{link}} \, \epsilon_f^2}}, \quad
g_{SU(3)} \sim \frac{1}{\sqrt{\rho_{\text{triplet}} \, \epsilon_f^2}}
\]

These expressions emerge purely from dimensional analysis and the statistical mechanics of bundle topology.

\section*{4. Derived Coupling Ratios}

Predicted ratios of Standard Model couplings:

\[
\frac{g_{SU(2)}}{g_{U(1)}} \sim \sqrt{\frac{\rho_{\text{wind}}}{\rho_{\text{link}}}}, \quad
\frac{g_{SU(3)}}{g_{SU(2)}} \sim \sqrt{\frac{\rho_{\text{link}}}{\rho_{\text{triplet}}}}
\]

These ratios are testable predictions of SAT and depend only on the relative abundance of topological structures in the filament network.

\section*{5. Physical Interpretation}

\begin{itemize}
    \item U(1) coupling strength emerges from winding loop prevalence.
    \item SU(2) from stable 2-filament links (e.g., twisted ribbons).
    \item SU(3) from triple-linking geometries (e.g., Borromean configurations).
    \item Couplings vary with spatial/temporal filament topology; cosmological or local deviations possible.
\end{itemize}

\section*{6. Falsifiability and Tests}

SAT predicts:
\begin{itemize}
    \item Ratios of gauge couplings must match linking density ratios within the filament network.
    \item Changes in observed coupling constants under extreme conditions (e.g., early universe) must trace to topology shifts.
    \item Discovery of stable 4-link bundles would falsify O4 and imply new couplings not captured by current invariants.
\end{itemize}

\section*{7. Module Linkages}

\begin{itemize}
    \item \textbf{O3}: Gauge symmetry structure originates in linking class algebra.
    \item \textbf{O6}: Unified action embeds couplings directly from statistical fields.
    \item \textbf{O8}: Couplings and mass suppression co-emerge from same topological network.
\end{itemize}

\section*{8. Frame Declaration}

This derivation is performed in \textbf{Interpretive Mode 1 (True Block)} with statistical topology as the sole source of dynamical coupling values.

\end{document}
